%% LUDUS Academik Journal starter — ludusofficial document class.
%%
%% This class file (ludusofficial.cls) IS bundled here: it is not part of any
%% TeX Live distribution, so a download without it would not compile.
%%
%% Class options: a colour theme (blue/red/green/purple/orange/cyan) and an
%% article type (fullpaper/letter/poster/survey/shortpaper/artpaper).
\documentclass[blue,fullpaper]{ludusofficial}

%% Document metadata — shown in the page header/footer banner.
\journalname{LUDUS}
\journalsubtitle{International Journal of Game Studies}
\conferencename{LUDUS 2026}
\publicationyear{2026}

%% DOI (optional — leave the line out entirely if not yet assigned)
%\articledoi{10.1234/ludus.2026.example}

%% Article information
\title{Difficulty Curves as Rhetorical Devices: Pacing in Roguelike Design}
\shorttitle{Difficulty Curves as Rhetoric}
\shortauthor{Ada Lovelace \& Grace Hopper}

\author{%
    \textbf{Ada Lovelace}\textsuperscript{1}, \textbf{Grace Hopper}\textsuperscript{2}\\[0.2cm]
    \mdseries\small\textsuperscript{1}Department of Game Design, University of Somewhere\\
    \small\textsuperscript{2}School of Interactive Media, Institute of Computing\\
    \small Corresponding: a.lovelace@example.ac.uk
}

\begin{document}

%% Abstract and keywords must be defined BEFORE \maketitle.
\begin{abstract}
Roguelikes are frequently described as "punishing," yet their difficulty
curves are rarely analysed as a deliberate communicative device rather than
an incidental design cost. We argue that permadeath-driven pacing functions
rhetorically: escalating failure states teach players to read a run's risk
profile before committing to it, turning death into feedback rather than
punishment. Drawing on case studies of three commercially released
roguelikes, we identify three pacing patterns — front-loaded attrition,
plateaued escalation, and spike-and-recover — and show that each
communicates a different relationship between player skill and player luck.
We conclude that difficulty curve design is best understood as an argument
the game makes about what kind of failure is meaningful.
\end{abstract}

\keywords{roguelike; difficulty curve; procedural rhetoric; game pacing; permadeath; player feedback}

%% Maketitle creates the full-width header with abstract and keywords, then
%% starts the two-column body layout.
\maketitle

%% From here, content is in two columns.
\section{Introduction}

Roguelikes make failure routine: a run ends, progress resets, and the player
begins again with only the knowledge gained, not the items or levels lost.
This design is often framed as punishing by contrast with genres that
preserve player progress unconditionally. We take a different view: a
roguelike's difficulty curve is not simply "how hard the game is" but a
structured argument about which kinds of failure the game treats as
meaningful information versus arbitrary bad luck.

\subsection{Why Pacing, Not Just Difficulty}

Difficulty alone is a scalar; pacing is a curve, and a curve has shape.
Two roguelikes can share an identical average failure rate while
communicating entirely different things about player agency, depending on
\emph{when} in a run failure becomes likely and how sharply the odds shift
once it does.

\section{Three Pacing Patterns}

We analysed run-length and cause-of-death data from three commercially
released roguelikes spanning distinct design lineages, and identify three
recurring pacing shapes.

\subsection{Front-Loaded Attrition}

Difficulty is highest in the opening minutes, when the player has the
fewest resources to respond to it, then eases as build-defining choices
compound. This pattern rewards players who treat early runs as
information-gathering rather than progress, echoing the "legible failure"
argument in procedural rhetoric~\cite{bogost2007}.

\subsection{Plateaued Escalation}

Difficulty rises in discrete steps rather than continuously, holding steady
across several encounters before jumping at fixed narrative or structural
checkpoints. Players can therefore read their current plateau as a signal
of what the game expects them to already know, similar to the
environmental legibility cues described by~\cite{jenkins2004}.

\subsection{Spike-and-Recover}

Isolated high-variance encounters interrupt an otherwise gentle curve,
after which difficulty drops back below the pre-spike baseline. This
pattern communicates that a single bad outcome is a local event, not
evidence the run itself is doomed — the closest of the three to Murray's
notion of a game world that remains legible and navigable even when hostile
to the player~\cite{murray1997}.

\begin{figure}[h!]
    \centering
    \includegraphics[width=\columnwidth]{testImg.jpg}
    \caption{Illustrative reference image used as a figure placeholder in this starter; replace with your own.}
    \label{fig:example}
\end{figure}

\section{Discussion}

Table~\ref{tab:patterns} summarises how each pattern differs in where
difficulty concentrates within a run and what it implies about the source
of failure.

\begin{table}[h!]
\centering
\caption{Three difficulty-curve pacing patterns and what each communicates about failure.}
\label{tab:patterns}
\begin{tabular}{lcc}
\hline
\textbf{Pattern} & \textbf{Peak timing} & \textbf{Implied cause} \\
\hline
Front-loaded  & Early    & Player unpreparedness \\
Plateaued     & Stepped  & Structural checkpoint \\
Spike-recover & Isolated & Local bad luck \\
\hline
\end{tabular}
\end{table}

None of these patterns is inherently more punishing in aggregate terms —
they can share an identical overall failure rate — but they differ sharply
in what a player is entitled to conclude from a death. A front-loaded curve
tells the player to treat early deaths as tuition; a spike-and-recover curve
tells them not to over-read a single bad encounter. Conflating these as
"just difficulty" erases a design decision that is doing real
communicative work.

\section{Conclusion}

Treating a roguelike's difficulty curve as a rhetorical structure rather
than a tuning parameter reframes what "fair" failure means in the genre: a
curve is fair not because it is gentle, but because it accurately signals
which failures are lessons and which are noise. Future work could extend
this framework beyond roguelikes to other permadeath-adjacent genres, and
test whether players' own accounts of "unfair" deaths track the pacing
pattern we identify here or diverge from it.

\begin{acknowledgments}
We thank playtesters who shared run logs used to identify the pacing
patterns discussed here.
\end{acknowledgments}

%% References using BibTeX.
%% Compile with: pdflatex -> bibtex -> pdflatex -> pdflatex (or xelatex in
%% place of pdflatex — this class supports both; see the Building note in
%% README.md for why xelatex is the default in engine.txt).
\bibliographystyle{plain}
\bibliography{references}

\end{document}
